\documentclass[11pt]{amsart}

\usepackage[T1]{fontenc}
\usepackage{lmodern}
\usepackage{microtype}
\usepackage{amsmath,amssymb,amsthm}
\usepackage[hidelinks]{hyperref}

\hypersetup{
  pdftitle={Counterexamples to the Chan--Pak--Panova Matching Conjecture}
}

\newtheorem{theorem}{Theorem}
\newtheorem{proposition}[theorem]{Proposition}

\newcommand{\cE}{\mathcal E}
\newcommand{\cN}{\mathcal N}

\title[Counterexamples to the Chan--Pak--Panova conjecture]
{Counterexamples to the Chan--Pak--Panova Matching Conjecture}

\date{}
\subjclass[2020]{06A07, 05C70, 05A20}
\keywords{linear extension, Stanley inequality, promotion, matching,
Hall's theorem}

\begin{document}

\begin{abstract}
Chan, Pak, and Panova proposed a promotion--demotion bipartite graph whose
conjectured matching would give a direct injection for Stanley's log-concavity
inequality. Using the inverse-order word sweeps that realize the stated
promotion--demotion effect, we disprove their Conjecture~9.4. For every
$m\ge2$, let $P_m$ have elements $x,d,b_1,\ldots,b_m$ and only nontrivial
comparability $x\prec d$, and take $a=2$. The resulting graph $G_m$ satisfies
\[
 |V|=m^2(m!)^2,\qquad
 |W|=(m^2-1)(m!)^2,\qquad
 \nu(G_m)=|\cN(W)|=(m-1)(m!)^2.
\]
Thus every maximum matching leaves $m(m-1)(m!)^2$ vertices of $W$ unmatched,
even though $|V|-|W|=(m!)^2>0$. The smallest counterexample has three
elements and is the antichain; the four-element poset $P_2$ is smallest among
counterexamples for which Stanley's inequality is strict.
\end{abstract}

\maketitle

\section{The matching conjecture}

Let $P=(X,\prec)$ be a finite poset, let $x\in X$, and let
$\cE(P,x,r)$ be the set of linear extensions placing $x$ in position $r$.
Stanley's inequality~\cite{Sta81} states that
\begin{equation}\label{eq:stanley}
 |\cE(P,x,a)|^2
 \ge
 |\cE(P,x,a-1)|\,|\cE(P,x,a+1)|.
\end{equation}
Chan, Pak, and Panova~\cite[Conjecture~9.4]{CPP23} proposed a bipartite graph
with
\[
 V=\cE(P,x,a)\times\cE(P,x,a),
 \qquad
 W=\cE(P,x,a-1)\times\cE(P,x,a+1),
\]
and conjectured that it has a matching saturating $W$.

We make the promotion--demotion convention explicit. Write a linear extension
as a word $w=w_1\cdots w_n$. For $1\le r<n$, let $\tau_r$ interchange
$w_r$ and $w_{r+1}$ when they are incomparable, and fix $w$ otherwise.
Section~7 of~\cite{CPP23} uses a right action, whereas Section~9.12 writes the
operators on the left. We use the inverse-order word sweeps that realize the
stated effect in Section~9.12:
\begin{equation}\label{eq:sweeps}
 D_i(w)=w\tau_{i-1}\cdots\tau_1,
 \qquad
 U_j(w)=w\tau_j\cdots\tau_{n-1}.
\end{equation}
Equivalently, these are induced by the left action $g\cdot w:=wg^{-1}$.
For $(\pi,\sigma)\in W$ and an element $y$ occurring after $x$ in $\pi$ and
before $x$ in $\sigma$, let $i$ and $j$ be the positions of $y$ in $\pi$ and
$\sigma$, respectively, and join $(\pi,\sigma)$ to
\[
 \bigl(D_i(\pi),U_j(\sigma)\bigr)\in V.
\]
For $S\subseteq W$, write $\cN(S)\subseteq V$ for its neighborhood, and let
$\nu(G)$ denote the maximum matching size.

\section{An infinite family}

\begin{theorem}\label{thm:main}
For every $m\ge2$, let
\[
 X_m=\{x,d,b_1,\ldots,b_m\}
\]
and let $P_m$ have only nontrivial comparability $x\prec d$. For the graph
$G_m$ above with $a=2$,
\begin{align*}
 |V|&=m^2(m!)^2,\\
 |W|&=(m^2-1)(m!)^2,\\
 \nu(G_m)=|\cN(W)|&=(m-1)(m!)^2.
\end{align*}
Consequently, under the convention in~\eqref{eq:sweeps},
Conjecture~9.4 of~\cite{CPP23} is false. Every maximum matching leaves
exactly
\[
 |W|-\nu(G_m)=m(m-1)(m!)^2
\]
vertices of $W$ unmatched, while
\[
 |V|-|W|=(m!)^2>0,
 \qquad
 \frac{|W|}{|\cN(W)|}=m+1.
\]
\end{theorem}

\begin{proof}
Set $B=\{b_1,\ldots,b_m\}$ and abbreviate
$\cE_r=\cE(P_m,x,r)$. Since $x\prec d$ is the only nontrivial
comparability,
\[
 |\cE_1|=(m+1)!,
 \qquad
 |\cE_2|=m\,m!,
 \qquad
 |\cE_3|=m(m-1)(m-1)!.
\]
Indeed, if $x$ is first, the remaining elements are arbitrary; if $x$ is
second, the first element is any member of $B$; and if $x$ is third, the
first two elements are an ordered pair of distinct members of $B$. Hence
\[
 |V|=|\cE_2|^2=m^2(m!)^2
\]
and
\[
 |W|=|\cE_1|\,|\cE_3|=(m^2-1)(m!)^2.
\]

For $y\in B$, let $\cN_y$ be the set of endpoints in $V$ of edges labeled by
$y$. The two entries before $x$ in every $\sigma\in\cE_3$ are distinct
members of $B$, so every edge label belongs to $B$. Since $y$ is incomparable
with every other element, each sweep in~\eqref{eq:sweeps} moves $y$ through
all intervening positions. It follows that $\cN_y$ consists exactly of the
pairs
\begin{equation}\label{eq:Ny}
 \bigl(yx\alpha,\,zx\beta y\bigr),
\end{equation}
where $z\in B\setminus\{y\}$, $\alpha$ is an arbitrary ordering of
$X_m\setminus\{x,y\}$, and $\beta$ is an arbitrary ordering of
$X_m\setminus\{x,y,z\}$. Conversely, every pair in~\eqref{eq:Ny} is obtained
from
\[
 \bigl(xy\alpha,\,zyx\beta\bigr)\in W
\]
by the edge labeled $y$. Therefore
\[
 |\cN_y|=m!\,(m-1)(m-1)!.
\]
The sets $\cN_y$ are pairwise disjoint, since $y$ is the first letter of the
first coordinate and the last letter of the second. Thus
\[
 |\cN(W)|
 =\sum_{y\in B}|\cN_y|
 =m\,m!\,(m-1)(m-1)!
 =(m-1)(m!)^2.
\]

Finally, the edges
\[
 \bigl(xy\alpha,zyx\beta\bigr)
 \;--\;
 \bigl(yx\alpha,zx\beta y\bigr)
\]
as $y,z,\alpha,\beta$ vary are pairwise disjoint on both sides and saturate
$\cN(W)$. Hence $\nu(G_m)=|\cN(W)|$. Since $|\cN(W)|<|W|$, Hall's condition
fails. The remaining formulas follow immediately.
\end{proof}

The strict inequality $|V|>|W|$ shows that the obstruction persists even
when Stanley's inequality~\eqref{eq:stanley} is strict. In particular, the
result refutes only the proposed matching construction, not Stanley's
inequality itself.

\section{Smallest counterexamples}

\begin{proposition}\label{prop:minimal}
Up to isomorphism preserving the distinguished element $x$, the antichain on
three elements is the unique smallest counterexample. Every three-element
instance with $W\ne\varnothing$ satisfies $|V|=|W|$. Consequently, $P_2$ is
smallest among counterexamples for which Stanley's inequality is strict.
\end{proposition}

\begin{proof}
For the antichain on $\{x,b,c\}$ with $a=2$,
\[
 \cE_1=\{xbc,xcb\},\qquad
 \cE_2=\{bxc,cxb\},\qquad
 \cE_3=\{bcx,cbx\}.
\]
Thus $|V|=|W|=4$, whereas
\[
 \cN(W)=\{(bxc,cxb),(cxb,bxc)\}.
\]
Hence no matching saturates $W$. No smaller instance is possible because
$1<a<n$ requires $n\ge3$.

Now let $|X|=3$ and suppose $W\ne\varnothing$. A linear extension placing $x$
first and one placing $x$ last imply that $x$ is incomparable with the other
two elements. Deleting or inserting $x$ gives bijections between the linear
extensions of $P\setminus\{x\}$ and each of $\cE_1,\cE_2,\cE_3$. Therefore
$|\cE_1|=|\cE_2|=|\cE_3|$, and hence $|V|=|W|$.

If the other two elements are comparable, write them as $b\prec c$. Then each
$\cE_r$ is a singleton, and the edge labeled $b$ joins $(xbc,bcx)$ to
$(bxc,bxc)$. Thus the antichain is the only three-element counterexample.
The poset $P_2$ has four elements and, by Theorem~\ref{thm:main}, satisfies
\[
 |V|=16,\qquad |W|=12,\qquad \nu(G_2)=4,
\]
so it is smallest in the strict regime.
\end{proof}

\begin{thebibliography}{9}

\bibitem{CPP23}
S.~H. Chan, I. Pak, and G. Panova,
\emph{Effective poset inequalities},
SIAM J. Discrete Math. \textbf{37} (2023), no.~3, 1842--1880.

\bibitem{Sta81}
R.~P. Stanley,
\emph{Two combinatorial applications of the Aleksandrov--Fenchel inequalities},
J. Combin. Theory Ser. A \textbf{31} (1981), no.~1, 56--65.

\end{thebibliography}

\end{document}